\documentclass[11pt]{article}

\usepackage[a4paper,margin=25mm]{geometry}
\usepackage{amsmath,amssymb,bm}
\usepackage{booktabs}
\usepackage{hyperref}
\usepackage{microtype}

\hypersetup{colorlinks=true,linkcolor=blue,urlcolor=blue,citecolor=blue}

\newcommand{\SO}{\mathrm{SO}(3)}
\newcommand{\R}{\mathbb{R}}
\newcommand{\Exp}{\operatorname{Exp}}
\newcommand{\Log}{\operatorname{Log}}
\newcommand{\diag}{\operatorname{diag}}
\newcommand{\transpose}{\mathsf{T}}

\title{Factor-Graph Model Predictive Control and\\
Terminal-Measurement Factor MPC}
\author{IPN\_MPC}
\date{\today}

\begin{document}
\maketitle

\begin{abstract}
This note introduces model predictive control (MPC) expressed as nonlinear factor-graph
optimization and a reduced-variable variant based on the
\texttt{TerminalAccelerationGyroMeasurementFactor}.  The conventional formulation represents
future states and controls as optimization variables.  The terminal-measurement formulation
instead rolls the current state forward through prefixes of the control sequence and compares each
prediction directly with a fixed trajectory set point.  The latter removes future-state variables,
uses analytical chain-rule Jacobians, employs a discrete Riccati terminal weight for local
stability, and enforces hard control bounds with a constrained quadratic-programming refinement.
\end{abstract}

\section{MPC as factor-graph optimization}

For a horizon of length $N$, let $x_k$ denote the state, $u_k$ the control input, and
$x_k^{\mathrm{ref}}$ the reference state.  A standard discrete-time MPC problem is
\begin{align}
\min_{x_{1:N},u_{0:N-1}}\quad
& \sum_{k=1}^{N-1}\frac{1}{2}
  \lVert x_k-x_k^{\mathrm{ref}}\rVert_{Q_k}^{2}
  +\sum_{k=0}^{N-1}\frac{1}{2}\lVert u_k-u_k^{\mathrm{ref}}\rVert_{R_k}^{2}
  +\frac{1}{2}\lVert x_N-x_N^{\mathrm{ref}}\rVert_{P}^{2},
\label{eq:standard-mpc-cost}\\
\text{subject to}\quad
&x_{k+1}=f(x_k,u_k),\qquad k=0,\ldots,N-1,\label{eq:dynamics}\
&u_{\min}\leq u_k\leq u_{\max}.\label{eq:bounds}
\end{align}
Here $\lVert e\rVert_W^2=e^\transpose W e$.  In a factor graph, every quadratic term is a
factor and every unknown state or control is a variable node.  A Gaussian factor with covariance
$\Sigma$ contributes
\begin{equation}
\frac{1}{2}e^\transpose\Sigma^{-1}e,
\end{equation}
so the desired weight matrix is supplied to GTSAM as the information matrix
$W=\Sigma^{-1}$.

Typical graph components are:
\begin{itemize}
  \item a strong prior on the measured current state $x_0$;
  \item dynamics factors joining adjacent states $x_k,u_k,x_{k+1}$;
  \item tracking factors on $x_k$;
  \item control regularization and smoothness factors on $u_k$ and $(u_{k-1},u_k)$;
  \item a terminal tracking factor on $x_N$; and
  \item optional obstacle, safety, and actuator constraints.
\end{itemize}

After optimization, MPC applies only $u_0$, advances the simulator or vehicle, shifts the previous
solution to form a warm start, and solves the horizon again.

\section{Acceleration--angular-speed model}

The terminal-factor applications use
\begin{equation}
x_k=(R_k,p_k,v_k),\qquad R_k\in\SO,quad p_k,v_k\in\R^3,
\end{equation}
and the six-dimensional control
\begin{equation}
u_k=\begin{bmatrix}a_k\\\omega_k\end{bmatrix}\in\R^6,
\end{equation}
where $a_k$ is world-frame acceleration and $\omega_k$ is body angular speed.  With sampling time
$\Delta t$, the rollout is
\begin{align}
p_{k+1} &= p_k+v_k\Delta t+\frac{1}{2}a_k\Delta t^2,\label{eq:position}\
v_{k+1} &= v_k+a_k\Delta t,\label{eq:velocity}\
R_{k+1} &= R_k\Exp(\omega_k\Delta t).\label{eq:rotation}
\end{align}
The optimized kinematic command is converted to collective thrust and torque before it is applied
to the full quadrotor simulator.

\section{State-variable terminal rollout factor}

The \texttt{TerminalAccelerationGyroFactor} connects the initial state, a control prefix, and an
optimized future state:
\begin{equation}
\{X_0,V_0,U_0,\ldots,U_{j-1},X_j,V_j\}.
\end{equation}
Starting from $x_0$, define
\begin{equation}
\widehat{x}_j=f^{(j)}(x_0,u_0,\ldots,u_{j-1}).
\end{equation}
Its residual is
\begin{equation}
e_j=
\begin{bmatrix}
p_j-\widehat p_j\\
\Log(\widehat R_j^{-1}R_j)\\
v_j-\widehat v_j
\end{bmatrix}\in\R^9,
\qquad
E_j=\frac{1}{2}e_j^\transpose Q_j e_j.
\label{eq:state-terminal-residual}
\end{equation}
The factor is a soft multiple-shooting constraint.  Separate priors or tracking factors attract
$X_j,V_j$ to the desired trajectory.  This formulation retains explicit predicted-state variables,
which are useful when other factors must attach directly to future states, for example obstacle,
sensor, or state-boundary factors.

\section{Terminal-measurement factor MPC}

\subsection{Fixed set point instead of a future-state variable}

The \texttt{TerminalAccelerationGyroMeasurementFactor} stores a desired pose and velocity as fixed
factor measurements.  Its variable set is only
\begin{equation}
\{X_0,V_0,U_0,\ldots,U_{j-1}\};
\end{equation}
there is no $X_j$ or $V_j$ variable.  For the measured set point
$z_j=(R_j^{\mathrm{ref}},p_j^{\mathrm{ref}},v_j^{\mathrm{ref}})$, the factor residual is
\begin{equation}
e_j=
\begin{bmatrix}
p_j^{\mathrm{ref}}-\widehat p_j\\
\Log(\widehat R_j^{-1}R_j^{\mathrm{ref}})\\
v_j^{\mathrm{ref}}-\widehat v_j
\end{bmatrix},
\qquad
E_j=\frac{1}{2}e_j^\transpose Q_j e_j.
\label{eq:measurement-residual}
\end{equation}
The MPC graph adds $N$ prefix factors:
\begin{align}
e_1 &= z_1-f(x_0,u_0),\nonumber\\
e_2 &= z_2-f^{(2)}(x_0,u_0,u_1),\nonumber\\
&\ \vdots\nonumber\\
e_N &= z_N-f^{(N)}(x_0,u_0,\ldots,u_{N-1}).
\end{align}
Consequently, the unknowns are the current state and control sequence, rather than an entire
state--control trajectory.

\subsection{Comparison of the two formulations}

\begin{table}[h]
\centering
\begin{tabular}{@{}lll@{}}
\toprule
Property & State-variable factor & Measurement factor \\
\midrule
Future $X_j,V_j$ variables & Required & Removed \\
Set point & Separate prior/factor & Stored in factor \\
Attach factors to future state & Direct & Requires rollout-dependent factor \\
Graph size & Larger & Smaller \\
Tracking residual & Future state minus rollout & Set point minus rollout \\
Application & \texttt{terminal\_acceleration\_gyro\_mpc} &
\texttt{terminal\_acceleration\_gyro\_setpoint\_mpc} \\
\bottomrule
\end{tabular}
\caption{Comparison of terminal acceleration--gyro MPC formulations.}
\end{table}

\section{Analytical chain-rule Jacobians}

Numerically differentiating every prefix repeats many rollouts and is unnecessarily expensive.
The measurement factor therefore uses analytical derivatives.  For a $j$-step rollout,
\begin{align}
\widehat v_j &= v_0+\Delta t\sum_{k=0}^{j-1}a_k,\\
\widehat p_j &= p_0+j\Delta t\,v_0
 +\Delta t^2\sum_{k=0}^{j-1}\left(j-k-\frac{1}{2}\right)a_k.
\end{align}
Thus,
\begin{align}
\frac{\partial e_p}{\partial v_0}&=-j\Delta t I,&
\frac{\partial e_v}{\partial v_0}&=-I,\\
\frac{\partial e_p}{\partial a_k}&=-\left(j-k-\frac{1}{2}\right)\Delta t^2I,&
\frac{\partial e_v}{\partial a_k}&=-\Delta t I.
\end{align}

For rotation, write
\begin{equation}
\widehat R_j=R_0E_0E_1\cdots E_{j-1},
\qquad E_k=\Exp(\omega_k\Delta t).
\end{equation}
Let $S_k=E_{k+1}\cdots E_{j-1}$ be the rotation suffix after $E_k$.  A local perturbation of
$E_k$ is transported to the final rotation through $\operatorname{Ad}_{S_k^{-1}}$.  Combining
this transport with the exponential-map Jacobian and the Jacobian of
$\Log(\widehat R_j^{-1}R_j^{\mathrm{ref}})$ gives
\begin{equation}
\frac{\partial e_R}{\partial\omega_k}
=J_{\Log}\,J_{\mathrm{between}}\,
\operatorname{Ad}_{S_k^{-1}}J_{\Exp}(\omega_k\Delta t)\Delta t.
\end{equation}
All suffixes are accumulated in one backward pass.  This avoids finite differences and keeps the
default 20-step MPC comfortably below its 10-ms solve-time target on the development machine.

\section{Terminal weight for local stability}

Near zero tracking error, use the local vector
\begin{equation}
\xi_k=\begin{bmatrix}\delta p_k&\delta\theta_k&\delta v_k\end{bmatrix}^\transpose
\in\R^9,
\qquad
\nu_k=\begin{bmatrix}\delta a_k&\delta\omega_k\end{bmatrix}^\transpose\in\R^6.
\end{equation}
The linearized model is
\begin{equation}
\xi_{k+1}=A\xi_k+B\nu_k,
\end{equation}
with
\begin{equation}
A=\begin{bmatrix}
I&0&\Delta t I\\
0&I&0\\
0&0&I
\end{bmatrix},
\qquad
B=\begin{bmatrix}
\tfrac{1}{2}\Delta t^2I&0\\
0&\Delta t I\\
\Delta t I&0
\end{bmatrix}.
\end{equation}
Given positive-definite stage weights $Q$ and $R$, the terminal information matrix $P$ is the
stabilizing solution of the discrete algebraic Riccati equation
\begin{equation}
P=Q+A^\transpose PA
-A^\transpose PB(R+B^\transpose PB)^{-1}B^\transpose PA.
\label{eq:dare}
\end{equation}
The feedback gain is
\begin{equation}
K=(R+B^\transpose PB)^{-1}B^\transpose PA.
\end{equation}
The implementation verifies that $\rho(A-BK)<1$, that the normalized Riccati residual is small,
and that the Lyapunov terminal cost strictly decreases under the local feedback.  This establishes
local stability of the unconstrained linearized error dynamics.  A full constrained nonlinear MPC
proof additionally requires a feasible invariant terminal set.

\section{Hard control-input bounds}

Standard GTSAM nonlinear optimizers do not directly impose inequality constraints.  The set-point
application first computes a nonlinear Levenberg--Marquardt solution and then applies a
sequential-quadratic-programming refinement using
\texttt{gtsam\_unstable::QPSolver}.  At an iterate $u_k$, the QP increment $\delta u_k$ obeys
\begin{align}
 \delta u_{k,i} &\leq u_{\max,i}-u_{k,i},\\
-\delta u_{k,i} &\leq u_{k,i}-u_{\min,i}.
\end{align}
These are exact box constraints for the linear control variables.  The QP is initialized with the
projected feasible increment
\begin{equation}
\delta u_k^{(0)}=\operatorname{clip}(u_k,u_{\min},u_{\max})-u_k.
\end{equation}
Before $u_0$ is applied, the complete optimized horizon is checked against the configured bounds
with a tolerance of $10^{-8}$.

\section{Receding-horizon algorithm}

At each control cycle, the set-point application performs the following steps:
\begin{enumerate}
  \item Read the current simulated pose and velocity and add strong priors on $X_0,V_0$.
  \item Generate $N$ circle-trajectory set points.
  \item Add one terminal-measurement factor for every control prefix.
  \item Add control regularization and control-smoothness factors.
  \item Use the DARE matrix $P$ for the final factor and $Q$ for intermediate factors.
  \item Solve the nonlinear least-squares problem from the shifted warm start.
  \item Refine the solution with hard box-constrained QPs.
  \item Verify feasibility and apply only $u_0$.
  \item Shift the control sequence and repeat at the next measurement.
\end{enumerate}

\section{Configuration and execution}

The principal settings are stored in
\texttt{config/terminal\_acceleration\_gyro\_mpc.yaml}.  The control ordering is
\begin{equation}
U(k)=[a_x,a_y,a_z,\omega_x,\omega_y,\omega_z]^\transpose.
\end{equation}
For example,
\begin{verbatim}
horizon: 20
dt: 0.10
terminal_position_sigma: 0.02
terminal_rotation_sigma: 0.02
terminal_velocity_sigma: 0.05
control_min: [-4.0, -4.0, -4.0, -2.0, -2.0, -2.0]
control_max: [ 4.0,  4.0,  4.0,  2.0,  2.0,  2.0]
hard_control_constraint_iterations: 2
\end{verbatim}

Build and run the two formulations with
\begin{verbatim}
cmake -S . -B build -DCMAKE_BUILD_TYPE=Release
cmake --build build --parallel

./build/terminal_acceleration_gyro_mpc \
  config/terminal_acceleration_gyro_mpc.yaml

./build/terminal_acceleration_gyro_setpoint_mpc \
  config/terminal_acceleration_gyro_mpc.yaml
\end{verbatim}
Append \texttt{--headless} to either command for automated testing without Pangolin visualization.

\section{Implementation map}

\begin{center}
\begin{tabular}{@{}ll@{}}
\toprule
Purpose & Source \\
\midrule
Terminal factor declarations &
\texttt{include/ipn\_mpc/control/dynamics\_control\_factor.h} \\
Factor residuals and Jacobians &
\texttt{src/control/dynamics\_control\_factor.cpp} \\
DARE terminal weight &
\texttt{src/control/lqr\_terminal\_weight.cpp} \\
State-variable MPC &
\texttt{apps/terminal\_acceleration\_gyro\_mpc.cpp} \\
Set-point MPC &
\texttt{apps/terminal\_acceleration\_gyro\_setpoint\_mpc.cpp} \\
Factor tests &
\texttt{tests/terminal\_dynamics\_factor\_test.cpp} \\
Terminal stability test &
\texttt{tests/lqr\_terminal\_weight\_test.cpp} \\
\bottomrule
\end{tabular}
\end{center}

\end{document}
